\subsection{Theoretical Implications}

The model has implications for three bodies of theory, all of them conditional on the propositions receiving empirical support.

For psychological contract theory, the main implication is that the formation context of a promise may shape what happens when it is broken. Formation research has emphasized schemas, promises, and the many channels through which obligations are conveyed \citep{rousseau2001schema}, and breach research has emphasized how employees interpret unmet obligations \citep{morrison1997employees,robinson2000development}. Promise publicity connects the two. If Proposition 1 holds, the public or private origin of an obligation belongs among the factors that explain why the same breach produces stronger violation in some employees than in others. It would also give ideological contract research, which has so far examined responses directed at the cause or the work \citep{yang2022going,deng2023serving}, a route into responses directed at the organization's public standing.

For voice and whistleblowing research, the model offers a way to bring the two streams into one frame. Voice research has mostly studied internal, upward communication \citep{morrison2014employee,morrison2023employee}, and whistleblowing research has treated disclosure to outside parties as one of several possible channels \citep{miceli2008whistleblowing,culiberg2017evolution}. Treating internal and public ethical voice as competing outcomes of the same felt violation, with internal voice safety deciding the balance, makes venue the thing to be explained. If Propositions 2 and 3 hold, an employee's choice to go public says as much about the promise and the internal channel as about the employee.

For generational research, the contribution is methodological as much as substantive. By stating cohort and career-stage accounts as rival propositions, the model gives researchers a way to test a generational claim without assuming it. This responds to repeated calls to separate cohort effects from age, period, and developmental processes \citep{parry2011generational,rudolph2017considering,rudolph2021generations}. It also connects to the view that generational identity becomes salient only under some conditions \citep{joshi2010unpacking}: if the cohort account receives support, the next question is whether it holds only when the breached promise was public, which is exactly where cohort-specific channel habits would be expected to matter.

\subsection{Practical Implications}

The practical implications below follow from the propositions and should be read as hypotheses for practice, not as evidence-based recommendations.

First, organizations that advertise ethical commitments to recruits are, on this account, creating promises that recruits may hold them to in public. Research on decoupling shows that ethics structures can be adopted without being built into everyday practice \citep{weaver1999integrated} and that adopted practices can fail to serve the ends they were meant to serve \citep{bromley2012smoke}. The model implies that the cost of such drift rises with the publicity of the commitment, because public promises are proposed to produce both stronger violation and a stronger pull toward public response. The implied remedy is not to say less but to make public ethical claims only where practice can support them.

Second, the internal channel has to be credible to the people least likely to trust it. If the career-stage account is right, junior employees are the group whose internal voice safety most needs attention. Managerial openness has been associated with voice through its effect on psychological safety \citep{detert2007leadership}, and socialization research shows that newcomers are still building role clarity and social acceptance \citep{bauer2007newcomer}. Extending that logic, onboarding that explains how ethical concerns are raised and what happens afterwards might raise internal voice safety early; whether it does is an open empirical question.

Third, the model suggests caution about responding to public voice mainly by restricting it. Research has shown that organizations' social media strategies and policies affect how sensitive employees are to the risks of whistleblowing through social media \citep{xiao2022employee}. Our model adds a prediction, not a finding: if restrictions reduce the attractiveness of public voice without raising internal voice safety, Proposition 4 implies that felt violation will be redirected toward exit and silence, which remove the concern from view without resolving it.

\subsection{A Testing Agenda}

Each proposition can be tested with existing research designs, although the constructs require some development.

Measurement is the first task. Breach and violation can be assessed with established measures developed in longitudinal psychological contract research \citep{robinson2000development}, with items adapted to the ethical content of ideological commitments as defined by \citet{thompson2003violations}. Internal voice safety can draw on the psychological safety measure of \citet{edmondson1999psychological}, reworded to refer to raising ethical concerns, and on measures of implicit voice theories \citep{detert2011implicit}. Internal ethical voice is close to prohibitive voice, for which a validated measure exists \citep{liang2012psychological}. The categories of intended response used in research on ethical culture, which range from inaction to reporting to management and external whistleblowing \citep{kaptein2011inaction}, provide a starting point for measuring the four responses together. Tests of the first two propositions should also account for organizational visibility. Large, consumer-facing organizations are likely both to make more public ethical claims and to draw more attention when criticized, which could produce an association between promise publicity and public voice for reasons unrelated to venue congruence. Promise publicity has no existing measure. Developing one would require items that ask employees where they learned about the employer's ethical commitments and how much weight each source carried, followed by the usual work on dimensionality and validity.

Two designs seem especially useful. The first is a multi-wave field study of organizational newcomers that deliberately includes both early-career entrants and older entrants, such as career changers. Measuring promise publicity at entry and breach, violation, internal voice safety, and responses at later waves would allow tests of Propositions 1 to 4 over time, following the logic of earlier longitudinal work on breach among new hires \citep{robinson2000development}. Including older newcomers is what makes Proposition 5 testable, because it lets cohort vary while tenure and position are held constant. A single study of this kind still cannot separate cohort from age, since the two move together at any one time; separating them fully would require repeating the design with later cohorts, a limitation that critics of generational research have stressed \citep{parry2011generational,rudolph2021generations}.

The second is a scenario experiment in which participants read about a breached ethical commitment and the experimenter varies whether the commitment was learned from the employer's public campaign or from a private conversation with a recruiter, crossed with cues about how safe internal reporting is. The wording of the commitment should be identical across conditions, so that publicity is not confounded with how specific the promise is. Such a design supports causal inference about Propositions 2 and 3. Its weakness is that it measures intentions, and meta-analytic evidence shows that the correlates of whistleblowing intentions do not carry over fully to actual whistleblowing \citep{mesmermagnus2005whistleblowing}. Field data on actual responses would therefore still be needed.

A third, qualitative design would address a part of the model that surveys and experiments handle poorly. Proposition 2 rests on a claim about how employees reason: that a promise made to the public is experienced as something the public is entitled to hear about when it is broken. Interviews with employees who have raised ethical concerns publicly, and with others who raised similar concerns internally, could show whether people describe their choice in these terms or in quite different ones, such as frustration with internal processes or a wish to warn prospective recruits. Evidence of that kind would refine the mechanism before larger quantitative tests are run.

Analytically, the four responses should not be forced into a single choice, because they can co-occur: an employee may raise a concern internally and later go public, or speak up while also looking for another job. Because the propositions concern relative propensities, each response is best measured separately, and the analysis should model their relative likelihood, for example by comparing an employee's propensity for public voice against the same employee's propensity for internal voice, rather than asking respondents to pick one response. Because Propositions 1 to 4 involve moderation and Proposition 5 involves comparisons across groups, sample sizes will need to be planned with interaction effects in mind.

\subsection{Boundary Conditions and Ways the Model Could Fail}

Several conditions could weaken or reverse the proposed relationships. Where public criticism of an employer carries serious legal or social penalties, the pull toward public voice should be suppressed regardless of promise publicity, and the model would predict more exit and silence instead. Where an organization makes no public ethical claims, promise publicity is close to its minimum and the model collapses into a standard breach account. The anonymity available on some platforms may lower the personal cost of public voice and so weaken the moderating role of internal voice safety. Retaliation is a further consideration: external whistleblowers have been found to suffer more extensive retaliation than internal ones \citep{dworkin1998internal}, and retaliation appears to depend on contextual factors \citep{mesmermagnus2005whistleblowing}, so anticipated retaliation may dampen public voice in ways the model does not capture.

The model could also simply be wrong. If public voice after an ethical breach turns out to depend mainly on the severity or the type of wrongdoing, with promise publicity adding nothing once those are controlled, Proposition 2 fails. If employees go public at similar rates whether the promise was public or private, venue congruence is not a useful idea. If internal voice safety does not moderate the venue of ethical voice, the model's account of why some employees stay inside would need another explanation, perhaps ethical culture at the organizational level \citep{kaptein2011inaction}. And if Proposition 5b receives consistent support, the Generation Z label would add nothing to the explanation; we regard that as a meaningful result, not a failure, since it would redirect attention to how organizations treat newcomers.

\subsection{Limitations}

The paper reports no data, and none of its propositions has been tested. Promise publicity is a new construct whose measurement and discriminant validity remain to be established. The model treats the choice of venue at a single point in time, whereas actual responses may unfold in sequence, for example with public voice following internal attempts that went nowhere; a process version of the model would capture this better. The literatures on which the model draws come mostly from North American and European settings, and the one source that directly describes a generational pattern in channel preferences is a non-peer-reviewed survey of stated intentions \citep{protect2025talkin}. The model's relevance elsewhere, and the pattern itself, remain open questions.
